%\exer{Application: déplacement d'un cavalier sur un échiquier}
%\vspace{-0.5cm}
\section*{Exercice d'application}


On donne la carte de la région Nouvelle Aquitaine.
\begin{figure}[H]
	\centering
		\includegraphics[width=.7\linewidth]{fig_01}
	\caption{Carte et graphe de la région Nouvelle-Aquitaine}
\end{figure}


\subsection*{Dictionnaire}
On souhaite créer un dictionnaire permettant de faire une correspondance entre le numéro des départements et leur nom.

\question{Implémenter une fonction \lstinline|add_dpt(d:{}, dpt:int, nom:str -> None| qui ajoute un département au dictionnaire  \lstinline{d}.}

\begin{exemple}[Tests]
\begin{lstlisting}
>>> dico = {}
>>> add_dpt(dico,64,"Pyrénées-Atlantiques")
>>> print(dico)
        {64:"Pyrénées-Atlantiques"}
\end{lstlisting}
\end{exemple}


\subsection*{Dictionnaire d'adjacence}
On souhaite modéliser la carte par un dictionnaire d'adjacence où les  sommets sont des numéros de départements et les valeurs sont la liste des départements limitrophes. 

\question{Donner l'instruction permettant de créer un graphe avec les Pyrénées-Atlantiques (64) et les Landes uniquement (40).}

\question{Implémenter la fonction \lstinline|add_dpt_d(G:{}, dpt:int, voisins : [int]) -> None| qui ajoute un département au graphe \lstinline{G}.}

\begin{exemple}[Tests]
\begin{lstlisting}
>>> aqui = {}
>>> add_dpt_d(aqui,64,[40])
>>> print(aqui)
        {64:[40], 40:[64]}
>>> add_dpt_d(aqui,47,[40,33,24])
>>> print(aqui)
        {64:[40], 40:[64,47], 47:[40,33,24]}
\end{lstlisting}
\end{exemple}

\question{Implémenter la fonction \lstinline|add_arete_d(G:{}, dpt1:int, dpt2:int) -> None| qui ajoute une arête entre deux départements.}


\question{Implémenter la fonction \lstinline|del_arete_d(G:{}, dpt:int) -> None| qui supprime un département de la région.}


\question{Implémenter la fonction \lstinline|degre_d(G:{}, dpt:int) -> int| qui donne le degré d'un sommet.}

\question{Implémenter la fonction \lstinline|degre_max_d(G:{}) -> int| qui renvoie le sommet de plus haut degré.}

\question{Implémenter la fonction \lstinline|degre_max_d(G:{}, d:{}) -> int| prend en argument le graphe de la région ainsi que le dictionnaire des départements qui renvoie le nom du (ou des) départements ayant le plus de départements limitrophes.}



\subsection*{Liste d'adjacence}
On souhaite modéliser la carte par une liste d'adjacence où les sommets sont des numéros de départements et les valeurs sont la liste des départements limitrophes.
\begin{marginfigure}
\includegraphics[width=\linewidth]{fig_02}
\end{marginfigure}
La numérotation des départements est modifiée conformément à la figure ci-contre.



\question{Donner l'instruction permettant de créer un graphe avec les départements 0, 1, 2, 3.}


\question{Implémenter la fonction \lstinline|add_dpt_l(G:[], dpt:int, voisins : [int]) -> None| qui ajoute un département au graphe \lstinline{G}.}

\begin{exemple}[Tests]
\begin{lstlisting}


\end{lstlisting}
\end{exemple}

\question{Implémenter la fonction \lstinline|add_arete_l(G:[], dpt1:int, dpt2:int) -> None| qui ajoute une arête entre deux départements.}

\question{Implémenter la fonction \lstinline|del_arete_l(G:[], dpt:int) -> None| qui supprime un département de la région.}

\question{Implémenter la fonction \lstinline|degre_d_l(G[], dpt:int) -> int| qui donne le degré d'un sommet.}

\question{Implémenter la fonction \lstinline|degre_max_d(G:[]) -> int| qui renvoie le sommet de plus haut degré.}


\subsection*{Matrice d'adjacence}
On souhaite modéliser la carte par une matrice d'adjacence. On conserve la numérotation précédente.


\question{Donner l'instruction permettant de créer un graphe avec les départements 0, 1, 2, 3.}


\question{Implémenter la fonction \lstinline|add_dpt_m(G:[[]], dpt:int, voisins : [int]) -> None| qui ajoute un département au graphe \lstinline{G}.}

\begin{exemple}[Tests]
\begin{lstlisting}


\end{lstlisting}
\end{exemple}

\question{Implémenter la fonction \lstinline|add_arete_m(G:[[]], dpt1:int, dpt2:int) -> None| qui ajoute une arête entre deux départements.}

\question{Implémenter la fonction \lstinline|del_arete_m(G:[[]], dpt:int) -> None| qui supprime un département de la région.}

\question{Implémenter la fonction \lstinline|degre_d_m(G[[]], dpt:int) -> int| qui donne le degré d'un sommet.}

\question{Implémenter la fonction \lstinline|degre_max_m(G:[[}]) -> int| qui renvoie le sommet de plus haut degré.}

